% Compile with XeLaTeX.
\documentclass[11pt,a4paper]{article}
\usepackage{fontspec,polyglossia}
\setmainlanguage{ukrainian}
\setmainfont{DejaVu Serif}
\usepackage{amsmath,amssymb}
\usepackage[margin=22mm]{geometry}
\setlength{\parindent}{0pt}
\setlength{\parskip}{5pt}
\newcommand{\unit}[1]{\,\mathrm{#1}}
\newcommand{\dd}{\mathrm{d}}
\title{Електрохімічні та оптичні методи\\ аналізу складу\\[4pt]\large Лекція 11. Математичний конспект}
\author{Метрологія, технологічні вимірювання та прилади}
\date{}
\begin{document}
\raggedright
\maketitle

\section{Кондуктометричний аналіз}

\textbf{Провідність комірки та електропровідність середовища.}
Для омічного відгуку $G=I/U=1/R_{\rm el}$, де $I$ --- струм (А), $U$ --- напруга (В), $R_{\rm el}$ --- опір (Ом). За змінного збудження $G$ тут означає провідність у відповідній моделі, а не модуль довільного комплексного адмітансу.
\[
 \kappa=K_{\rm cell}G,\qquad
 K_{\rm cell}=\frac{L}{S}\quad\text{для ідеалізованого однорідного поля}.
\]
$G$ --- См, $\kappa$ --- См/м, $K_{\rm cell}$ --- м$^{-1}$, $L$ --- м, $S$ --- м$^2$. Для реального зонда $K_{\rm cell}$ визначають за опорною електропровідністю і придатною моделлю.

\textbf{Гранична адитивність іонних внесків.}
За нескінченного розбавлення молярна провідність електроліту
$\Lambda_{\rm m}^\circ=\sum_i\nu_i\lambda_i^\circ$,
де $\nu_i$ --- стехіометричні кількості іонів, $\lambda_i^\circ$ --- їхні граничні молярні провідності (См$\cdot$м$^2$/моль).
У розбавленій суміші $\kappa\approx\sum_i\lambda_i^\circ c_i$ з $c_i$ у моль/м$^3$. Одне скалярне значення $\kappa$ загалом не визначає всіх $c_i$.

\textbf{Локальна температурна поправка.}
За фіксованого складу та в заданому інтервалі
\[
 \kappa(T)\simeq\kappa(T_0)[1+\alpha(T-T_0)],
 \qquad
 \kappa(T_0)\simeq\frac{\kappa(T)}{1+\alpha(T-T_0)}.
\]
$T,T_0$ --- К, температурна різниця у К, $\alpha$ --- К$^{-1}$; знаменник має бути додатним. За $T=T_0$ поправка тотожна. Для $\kappa(T)=0.110$ См/м, $\alpha=0.02$ К$^{-1}$ і $T-T_0=5$ К відновлюємо $\kappa(T_0)=0.100$ См/м.

\textbf{Табличне обернення.}
Для однієї матриці й температури за строго монотонної характеристики і сусідніх різних вузлів $\kappa_i,\kappa_{i+1}$ кусково-лінійна оцінка
\[
 \widehat c=c_i+(c_{i+1}-c_i)
  \frac{\kappa-\kappa_i}{\kappa_{i+1}-\kappa_i}.
\]
Область: відрізок між вузлами, включно з кінцями. Дробовий множник безрозмірний. Пряме і зворотне лінійні перетворення на тому самому відрізку взаємно обернені.

\section{Потенціометрія і термінологія полярографії}

\textbf{Рівняння Нернста.}
Для заданої відновної напівреакції з $z>0$ електронами і безрозмірним реакційним відношенням $Q_r$:
\[
 E=E^\circ-\frac{RT}{zF}\ln Q_r .
\]
$E,E^\circ$ --- В, $T$ --- К, $R\approx8.31446$ Дж/(моль$\cdot$К), $F\approx96485$ Кл/моль. Якщо відгук за обраної полярності має форму $E=E_0+(RT/F)\ln a_{\rm H^+}$, то
\[
 \mathrm{pH}=-\log_{10}a_{\rm H^+},
 \qquad E=E_0-\frac{RT\ln 10}{F}\,\mathrm{pH}.
\]
Активність $a_{\rm H^+}$ безрозмірна; практичну шкалу реалізують за допомогою буферів. За $T=298.15$ К ідеальний нахил становить приблизно $-59.16$ мВ на одиницю pH для наведеної полярності. Компенсація нахилу електрода не визначає температурної зміни pH самої проби.

\textbf{Вхідне навантаження.}
В ідеалізованій резистивній моделі $E_{\rm obs}=E\,R_{\rm in}/(R_s+R_{\rm in})$, де опори $R_s,R_{\rm in}$ в Омах. За $R_{\rm in}/R_s\to\infty$ маємо $E_{\rm obs}\to E$.

\textbf{Двоточкова калібрувальна характеристика.}
Для різних буферних значень $p_1,p_2$ і потенціалів $E_1,E_2$ у мВ:
\[
 m=\frac{E_2-E_1}{p_2-p_1},\qquad b=E_1-mp_1,\qquad
 \widehat{\mathrm{pH}}=\frac{E-b}{m}.
\]
$m$ --- мВ на одиницю pH, $b$ --- мВ. Необхідно $m\ne0$, а його знак і величина мають відповідати вимірювальній схемі та її критеріям. У моделі інтерполяції $\widehat{\mathrm{pH}}$ має лежати між $p_1$ і $p_2$; підстановка обох опорних потенціалів повертає відповідні pH.

\textbf{Модель полярографічної хвилі.}
Реєструють $I(E)$ (А) за заданого $E$ (В). За фіксованих умов локальну характеристику чистого струму хвилі можна записати як $I_{\rm lim}-I_{\rm bg}=k_I c$ з $k_I$ у А/(одиниця $c$). Обернення вимагає $k_I\ne0$ і каліброваного інтервалу; перекриті внески не можна окремо відновити з одного сумарного струму.

\section{Радіоізотопний і термохімічний аналіз}

\textbf{Матрична залежність.}
Нехай $c$ --- концентрація або частка цільового компонента, $T$ --- температура (К), $\mathbf u$ --- параметри матриці й геометрії. Загальна модель $s=f(c,T,\mathbf u)$ не обов’язково однозначно обернена за $c$. Для однорідної суміші та вузького моноенергетичного пучка без внеску розсіяних квантів:
\[
 -\ln(I/I_0)=\rho\,\ell\sum_i w_i\mu_{{\rm m},i}.
\]
$w_i$ --- масові частки ($\sum_i w_i=1$), $\mu_{{\rm m},i}$ --- м$^2$/кг, $\rho$ --- кг/м$^3$, $\ell$ --- м, $I,I_0$ мають однакові одиниці. Сума є масовим коефіцієнтом ослаблення суміші у прийнятій моделі. Один відлік не розділяє $\rho$, $\ell$ та невідомі $w_i$.

\textbf{Теплова потужність реакції.}
В ідеалізованій стаціонарній моделі екзотермічної реакції $\dot Q_{\rm rel}=-\Delta_r H\,\dot\xi$, де $\Delta_r H<0$ --- Дж/моль, $\dot\xi$ --- швидкість перебігу реакції (моль/с), $\dot Q_{\rm rel}$ --- Вт. Для сталого теплового відведення $H_{\rm th}>0$ (Вт/К) додаткове нагрівання $\Delta T\simeq\dot Q_{\rm rel}/H_{\rm th}$. Відгук залежить від перебігу реакції і теплопередачі.

\textbf{Локальна калібрувальна пряма.}
За фіксованих $T,\mathbf u$ і в каліброваному інтервалі
\[
 s=s_0+kc,\qquad \widehat c=(s-s_0)/k,\quad k\ne0.
\]
Наприклад, $s,s_0$ --- мВ, $c$ --- об.\,\%, $k$ --- мВ/(об.\,\%). За $s=s_0$ отримуємо нульове $c$, якщо нуль входить до області чинності. Якщо умови змінилися, належність сигналу числовому діапазону не підтверджує чинність обернення.

\section{Оптичні методи та межі їх застосування}

\textbf{Пропускання, абсорбанс і закон Бугера--Ламберта--Бера.}
\[
 \tau_{\rm opt}=\frac{I}{I_0},\qquad
 A=-\log_{10}\tau_{\rm opt},\qquad
 A=\varepsilon(\lambda)\ell c,\qquad
 \tau_{\rm opt}=10^{-A}.
\]
$I,I_0$ --- оптичні потужності або пропорційні їм скориговані сигнали в однакових одиницях, $I_0>0$. $\ell$ --- см, $c$ --- моль/л, $\varepsilon$ --- л/(моль$\cdot$см), $\lambda$ --- довжина хвилі. $A$ і $\tau_{\rm opt}$ безрозмірні. Закон застосовний для однорідної проби, вузької спектральної смуги, сталого стану поглинача та при нехтовному розсіюванні й паразитному світлі. За $\tau_{\rm opt}=1$ маємо $A=0$; за $\tau_{\rm opt}=0.1$ маємо $A=1$. Логарифм вимагає $\tau_{\rm opt}>0$. За $\tau_{\rm opt}>1$ маємо $A<0$ відносно опори, але це поза моделлю невід’ємної концентрації з $\varepsilon\ell>0$.

\textbf{Паразитне світло.}
За однакового додаткового сигналу $I_s>0$ для проби й опори
$\tau_{\rm obs}=(I+I_s)/(I_0+I_s)$. Якщо $0<I<I_0$, то $\tau_{\rm obs}>\tau_{\rm opt}$ і $A_{\rm obs}<A$. Пропускання, бокове розсіювання та абсорбанс не можна ототожнювати; зміна розподілу частинок змінює залежність бокового сигналу від складу.

\textbf{Рефракція та оптичне обертання.}
Для межі призми $n_p$ і проби $n_s<n_p$ закон Снеліуса дає
$\sin\theta_{\rm cr}=n_s/n_p$ для критичного кута, відліченого від нормалі всередині призми. Обидва показники безрозмірні й залежать від $\lambda$, а $n_s$ також від температури та складу.
За одного оптично активного компонента у придатному режимі
$\theta=[\alpha]_\lambda^T\ell c_m$, де $\theta$ --- у градусах, $\ell$ --- дм, $c_m$ --- г/мл, а $[\alpha]_\lambda^T$ --- градус$\cdot$мл/(дм$\cdot$г).
Знак $\theta$ визначає напрямок відносно узгодженої конвенції. У лінійній моделі суміші внески компонентів додаються, тому нульовий кут не виключає ненульові концентрації.
\end{document}
